\documentclass[11pt, a4paper]{article}

\usepackage{lmodern}
\usepackage{microtype}
\usepackage[margin=2.5cm]{geometry}
\usepackage[dvipsnames, table]{xcolor}
\definecolor{linkblue}{rgb}{0.0, 0.2, 0.6}
\definecolor{linkred}{rgb}{0.7, 0.13, 0.13}
\usepackage[style=numeric-comp, backend=biber, sorting=none]{biblatex}
\usepackage{citemsp} %──────────────────────────────────────────────────────── new package
\usepackage[colorlinks=true, citecolor=linkblue,
            linkcolor=linkred, urlcolor=linkblue]{hyperref}
\usepackage{tcolorbox}
\usepackage{authblk}
\usepackage{graphicx}
\usepackage{array}
\usepackage{booktabs}
\usepackage{orcidlink}
\usepackage{fontawesome5}
\usepackage{titlesec}

% ── Section formatting ────────────────────────────────────────────────────────
\titleformat{\section}
  {\Large\sffamily\bfseries\color{black!80}}
  {\colorbox{black!10}{\makebox[1.2cm]{\thesection}}}
  {1em}
  {}
  [\vspace{-0.3em}{\color{black!20}\rule{\textwidth}{0.5pt}}]

\titleformat{\subsection}
  {\large\sffamily\bfseries\color{black!70}}
  {\thesubsection}
  {1em}
  {}

% ── Helpers ───────────────────────────────────────────────────────────────────
\newcommand{\githublink}[1]{\href{https://github.com/#1}{\faGithub}}

% ── Example bibliography ─────────────────────────────────────────────────────
\begin{filecontents}[force]{refs-citemsp.bib}
@article{einstein1915,
  author  = {Einstein, Albert},
  title   = {Die Feldgleichungen der Gravitation},
  journal = {Sitzungsberichte der Preussischen Akademie der Wissenschaften},
  year    = {1915},
  pages   = {844--847}
}
@article{schwarzschild1916,
  author  = {Schwarzschild, Karl},
  title   = {{\"U}ber das Gravitationsfeld eines Massenpunktes nach der Einsteinschen Theorie},
  journal = {Sitzungsberichte der Preussischen Akademie der Wissenschaften},
  year    = {1916},
  pages   = {189--196}
}
@book{wald1984,
  author    = {Wald, Robert M.},
  title     = {General Relativity},
  publisher = {University of Chicago Press},
  year      = {1984}
}
@article{hawking1974,
  author  = {Hawking, Stephen W.},
  title   = {Black hole explosions?},
  journal = {Nature},
  volume  = {248},
  pages   = {30--31},
  year    = {1974}
}
@book{mtw1973,
  author    = {Misner, Charles W. and Thorne, Kip S. and Wheeler, John Archibald},
  title     = {Gravitation},
  publisher = {W. H. Freeman},
  year      = {1973}
}
@article{penrose1965,
  author  = {Penrose, Roger},
  title   = {Gravitational Collapse and Space-Time Singularities},
  journal = {Physical Review Letters},
  volume  = {14},
  pages   = {57--59},
  year    = {1965}
}
@article{kerr1963,
  author  = {Kerr, Roy P.},
  title   = {Gravitational Field of a Spinning Mass as an Example of Algebraically Special Metrics},
  journal = {Physical Review Letters},
  volume  = {11},
  pages   = {237--238},
  year    = {1963}
}
@article{chandrasekhar1931,
  author  = {Chandrasekhar, Subrahmanyan},
  title   = {The Maximum Mass of Ideal White Dwarfs},
  journal = {The Astrophysical Journal},
  volume  = {74},
  pages   = {81--82},
  year    = {1931}
}
\end{filecontents}
\addbibresource{refs-citemsp.bib}

% ── Title ─────────────────────────────────────────────────────────────────────
\title{\textbf{\texttt{citemsp}: Citation Locators for \LaTeX}}

\author[1]{Apostolos Tsampodimos%
  \thanks{\texttt{apostolos.tsampodimos@ftmc.lt}}%
  \,\,\githublink{ApostolosTsabodimos/citemsp-Latex-package}%
  \,\,\orcidlink{0009-0001-0230-5647}}
\author[2]{Nathaniel Sherrill%
  \thanks{\texttt{nathaniel.sherrill@itp.uni-hannover.de}}%
  \,\,\orcidlink{0000-0001-6909-3872}}
\author[2]{Agnese Mariotti%
  \thanks{\texttt{agnese.mariotti@itp.uni-hannover.de}}%
  \,\,\orcidlink{0000-0003-4602-7842}}

\affil[1]{Center for Physical Sciences and Technology,
  Saul\.etekio~3, 10257 Vilnius, Lithuania}
\affil[2]{Institute for Theoretical Physics, Leibniz University Hannover,
  Appelstra{\ss}e~2, 30167 Hannover, Germany}

\date{\today}

\begin{document}

\maketitle

\begin{abstract}
\noindent
We present \texttt{citemsp}, a \LaTeX{} package that adds per-key locators to numeric citation labels, rendered as compact superscript/subscript pairs\footnote{The package is available at \url{https://github.com/ApostolosTsabodimos/citemsp-Latex-package} and on CTAN.}. Working on top of either \texttt{biblatex} or \texttt{natbib}, the package provides a slash-delimited syntax with ten built-in prefix types --- covering sections, equations, theorems, figures, and more --- and a user-extensible registry for domain-specific locators.
\end{abstract}


\section{Syntax and Usage}

The \texttt{\textbackslash citemsp} command accepts a comma-separated list of citation keys with optional locators separated by forward slashes. By default, the first locator slot renders as \S\ (section) and the second as \P\ (paragraph). A single-letter prefix changes the locator type. With no locators, \texttt{\textbackslash citemsp} behaves identically to \texttt{\textbackslash cite}, including range compression when using \texttt{numeric-comp} or \texttt{sort\&compress}.

\begin{tcolorbox}[colback=gray!5, colframe=gray!40, boxrule=0.5pt, arc=2pt, top=6pt, bottom=6pt, left=10pt, right=10pt]
\begin{tabular*}{\dimexpr\linewidth-0pt}{@{}l@{\extracolsep{\fill}}c@{\quad}l@{}}
\texttt{\textbackslash citemsp\{key\}}                    & $\longrightarrow$ & \citemsp{wald1984}\\[3pt]
\texttt{\textbackslash citemsp\{key/section\}}             & $\longrightarrow$ & \citemsp{wald1984/6.2}\\[3pt]
\texttt{\textbackslash citemsp\{key/section/paragraph\}}   & $\longrightarrow$ & \citemsp{wald1984/6.2/3}\\[3pt]
\texttt{\textbackslash citemsp\{key//paragraph\}}          & $\longrightarrow$ & \citemsp{wald1984//3}\\[3pt]
\texttt{\textbackslash citemsp\{key/section/e4.3\}}        & $\longrightarrow$ & \citemsp{wald1984/6.2/e4.3}
\end{tabular*}
\end{tcolorbox}

\noindent Currently there exist ten built-in prefixes, which are listed in Table~\ref{tab:prefixes}. Each prefix can appear in either the superscript or subscript slot. Users can register additional prefixes with \texttt{\textbackslash citemspprefix\{letter\}\{label\}}.

\begin{table}[ht]
\centering
\rowcolors{2}{blue!3}{white}
\begin{tabular}{>{\raggedright\arraybackslash}p{0.46\textwidth}
                >{\raggedright\arraybackslash}p{0.46\textwidth}}
\toprule
\rowcolor{blue!15}
\textbf{Standard biblatex} & \textbf{citemsp} \\
\midrule
\verb|\autocite[§3.2]{smith2003}| \newline
$\rightarrow$ [1, \S3.2]
&
\verb|\citemsp{smith2003/3.2}| \newline
$\rightarrow$ \citemsp{smith2003/3.2}
\\[6pt]
\verb|\autocites[§3.2]{smith2003}| \verb|[eq.~4.3]{meyerrochow2003}| \newline
$\rightarrow$ [1, \S3.2][2, eq.~4.3]
&
\verb|\citemsp{smith2003/3.2,| \verb|meyerrochow2003/e4.3}| \newline
$\rightarrow$ \citemsp{smith2003/3.2, meyerrochow2003/e4.3}
\\[6pt]
\verb|\autocites[§3.2, ¶5]{smith2003}| \newline
$\rightarrow$ [1, \S3.2, \P5]
&
\verb|\citemsp{smith2003/3.2/5}| \newline
$\rightarrow$ \citemsp{smith2003/3.2/5}
\\[6pt]
No built-in typed locator syntax; user must type symbols manually
&
Ten built-in prefixes; extensible via \cmd{citemspprefix}
\\[6pt]
Range compression via \texttt{numeric-comp}
&
Plain entries (no~\texttt{/}) compressed via the backend's native
\cmd{cite}, so compression works identically
\\
\bottomrule
\end{tabular}
\caption{Feature comparison between standard \pkg{biblatex} locators and
  \pkg{citemsp}.}
\label{tab:comparison}
\end{table}


\section{Examples}

\subsection{Basic Usage}


\begin{tcolorbox}[colback=blue!3, colframe=blue!30, boxrule=0.5pt, arc=2pt, top=6pt, bottom=6pt, left=10pt, right=10pt]
\begin{verbatim}
The Schwarzschild metric~\citemsp{wald1984/6.1/2} describes
spherically symmetric vacuum solutions.
\end{verbatim}
\end{tcolorbox}

\noindent\textbf{Output:} The Schwarzschild metric~\citemsp{wald1984/6.1/2} describes spherically symmetric vacuum solutions.

\medskip
\noindent Multiple sources with independent locators:

\begin{tcolorbox}[colback=blue!3, colframe=blue!30, boxrule=0.5pt, arc=2pt, top=6pt, bottom=6pt, left=10pt, right=10pt]
\begin{verbatim}
Black hole thermodynamics~\citemsp{hawking1974/1/2,
wald1984/12.5/1} establishes a deep connection between
gravity and quantum mechanics.
\end{verbatim}
\end{tcolorbox}

\noindent\textbf{Output:} Black hole thermodynamics~\citemsp{hawking1974/1/2, wald1984/12.5/1} establishes a deep connection between gravity and quantum mechanics.

\medskip
\noindent Locator entries can be freely mixed:

\begin{tcolorbox}[colback=blue!3, colframe=blue!30, boxrule=0.5pt, arc=2pt, top=6pt, bottom=6pt, left=10pt, right=10pt]
\begin{verbatim}
General covariance~\citemsp{einstein1915, wald1984/4,
hawking1974/1/2} forms the foundation of general relativity.
\end{verbatim}
\end{tcolorbox}

\noindent\textbf{Output:} General covariance~\citemsp{einstein1915, wald1984/4, hawking1974/1/2} forms the foundation of general relativity.


\subsection{Similarity to {\textbackslash cite}}

With no locators, \texttt{\textbackslash citemsp} passes keys through the backend's native \texttt{\textbackslash cite}, producing identical output:

\medskip
\noindent
\verb|\cite{wald1984}| $\longrightarrow$ \cite{wald1984} \qquad
\verb|\citemsp{wald1984}| $\longrightarrow$ \citemsp{wald1984}

\medskip
\noindent
\begin{tabular}{@{}p{7cm}l@{}}
All plain (5 consecutive):     & \citemsp{einstein1915, schwarzschild1916, wald1984, hawking1974, mtw1973} \\[4pt]
All plain (non-consecutive):   & \citemsp{einstein1915, schwarzschild1916, wald1984, penrose1965, kerr1963, chandrasekhar1931} \\[4pt]
Range, then locator:           & \citemsp{einstein1915, schwarzschild1916, wald1984, hawking1974/3.2} \\[4pt]
Locator, then range:           & \citemsp{einstein1915/3.2, schwarzschild1916, wald1984, hawking1974, mtw1973} \\[4pt]
Range, locator, range:         & \citemsp{einstein1915, schwarzschild1916, wald1984, hawking1974/e2.1, mtw1973, penrose1965, kerr1963} \\[4pt]
Adjacent locators only:        & \citemsp{wald1984/3.1, hawking1974/e2.1, mtw1973/p42} \\
\end{tabular}


\subsection{Mixed Types and Custom Prefixes}

All built-in locator types can be freely mixed in a single call:

\begin{tcolorbox}[colback=blue!3, colframe=blue!30, boxrule=0.5pt, arc=2pt, top=6pt, bottom=6pt, left=10pt, right=10pt]
\begin{verbatim}
See~\citemsp{wald1984/T6.2, einstein1915/3.2/d5,
mtw1973/p520, hawking1974/1/e4.3} for further details.
\end{verbatim}
\end{tcolorbox}

\noindent\textbf{Output:} See~\citemsp{wald1984/T6.2, einstein1915/3.2/d5, mtw1973/p520, hawking1974/1/e4.3} for further details.

\medskip
\noindent Users can register domain-specific prefixes with \texttt{\textbackslash citemspprefix}. For example, a cosmology paper might define prefixes for window functions and sky patches:

\citemspprefix{w}{w.}
\citemspprefix{S}{\ensuremath{\Omega}}

\begin{tcolorbox}[colback=blue!3, colframe=blue!30, boxrule=0.5pt, arc=2pt, top=6pt, bottom=6pt, left=10pt, right=10pt]
\begin{verbatim}
\citemspprefix{w}{w.}
\citemspprefix{S}{\ensuremath{\Omega}}
The top-hat window~\citemsp{wald1984/4.2/w3} smooths
the power spectrum over sky patch~\citemsp{mtw1973/S12}.
\end{verbatim}
\end{tcolorbox}

\noindent\textbf{Output:} The top-hat window~\citemsp{wald1984/4.2/w3} smooths the power spectrum over sky patch~\citemsp{mtw1973/S12}.


\subsection{Paragraph example}

The following paragraph illustrates how \texttt{\textbackslash citemsp} reads in running text, mixing plain citations, section locators, and prefix types:

\begin{tcolorbox}[colback=blue!3, colframe=blue!30, boxrule=0.5pt, arc=2pt, top=6pt, bottom=6pt, left=10pt, right=10pt]
\begin{verbatim}
The Schwarzschild solution~\citemsp{schwarzschild1916/1/3}
was found shortly after the field
equations~\citemsp{einstein1915} were published.  A modern
treatment can be found in standard
textbooks~\citemsp{wald1984/6.1/2, mtw1973/p520}, where the
derivation proceeds via the Birkhoff
theorem~\citemsp{wald1984/T6.2}.  The singularity
theorems~\citemsp{penrose1965, hawking1974/1/e4.3} later
showed that gravitational collapse is generic.  For further
context see~\citemsp{kerr1963, chandrasekhar1931}.
\end{verbatim}
\end{tcolorbox}

\noindent\textbf{Output:} The Schwarzschild solution~\citemsp{schwarzschild1916/1/3} was found shortly after the field equations~\citemsp{einstein1915} were published.  A modern treatment can be found in standard textbooks~\citemsp{wald1984/6.1/2, mtw1973/p520}, where the derivation proceeds via the Birkhoff theorem~\citemsp{wald1984/T6.2}.  The singularity theorems~\citemsp{penrose1965, hawking1974/1/e4.3} later showed that gravitational collapse is generic.  For further context see~\citemsp{kerr1963, chandrasekhar1931}.


\section{Loading and Configuration}

\subsection*{With biblatex}

\begin{tcolorbox}[colback=gray!5, colframe=gray!40, boxrule=0.5pt, arc=2pt, top=6pt, bottom=6pt, left=10pt, right=10pt]
\begin{verbatim}
\usepackage[style=numeric-comp, backend=biber]{biblatex}
\addbibresource{refs.bib}
\usepackage{citemsp}
\end{verbatim}
\end{tcolorbox}

\noindent Use \texttt{numeric-comp} (or \texttt{numeric}) as the style. Range compression of consecutive plain entries requires \texttt{numeric-comp}.

\subsection*{Configuration}

\noindent The locator glyph size is controlled by \texttt{\textbackslash citemspscale} (default~0.35):

\medskip
\noindent
Default (0.35): \citemsp{wald1984/6.1/e4.3}
\qquad
{%
\renewcommand{\citemspscale}{0.25}%
Smaller (0.25): \citemsp{wald1984/6.1/e4.3}%
}
\qquad
{%
\renewcommand{\citemspscale}{0.50}%
Larger (0.50): \citemsp{wald1984/6.1/e4.3}%
}

\medskip
\noindent
To adjust the locator glyph size, redefine \texttt{\textbackslash citemspscale} after loading:

\begin{tcolorbox}[colback=gray!5, colframe=gray!40, boxrule=0.5pt, arc=2pt, top=6pt, bottom=6pt, left=10pt, right=10pt]
\begin{verbatim}
\usepackage{citemsp}
\renewcommand{\citemspscale}{0.4}  % default is 0.35
\end{verbatim}
\end{tcolorbox}


\section{Relation to Standard Locators}

The standard \texttt{biblatex} provides locators via its postnote mechanism.
The \texttt{\textbackslash citemsp} package replaces that mechanism with a
more compact notation. We illustrate the differences below.

\subsection*{Side-by-side source comparison}

Consider citing three sources, each with a different locator type:

\begin{tcolorbox}[colback=gray!5, colframe=gray!40, boxrule=0.5pt, arc=2pt, top=6pt, bottom=6pt, left=10pt, right=10pt]
\begin{verbatim}
% Standard biblatex (71 characters):
\autocites[§6.1, ¶2]{wald1984}[pp.\,520]{mtw1973}
  [§1, eq.\,4.3]{hawking1974}

% citemsp equivalent (53 characters):
\citemsp{wald1984/6.1/2, mtw1973/p520,
  hawking1974/1/e4.3}
\end{verbatim}
\end{tcolorbox}

\noindent\textbf{Rendered output:}

\medskip
\noindent
\begin{tabular}{@{}ll@{}}
biblatex: & \autocites[§6.1, ¶2]{wald1984}[pp.\,520]{mtw1973}[§1, eq.\,4.3]{hawking1974} \\[4pt]
citemsp:  & \citemsp{wald1984/6.1/2, mtw1973/p520, hawking1974/1/e4.3}
\end{tabular}

\medskip
\noindent The \texttt{citemsp} version is 25\% shorter in source characters
and uses a flat, comma-separated list rather than nested optional arguments.
Locators are rendered as superscript/subscript pairs attached to each
citation number, rather than appended as flat text inside the brackets.

\subsection*{Scaling to dense citations}

The difference becomes pronounced in review paragraphs with many
locator-bearing citations. Consider six sources:

\begin{tcolorbox}[colback=gray!5, colframe=gray!40, boxrule=0.5pt, arc=2pt, top=6pt, bottom=6pt, left=10pt, right=10pt]
\begin{verbatim}
% Standard biblatex (6 sources, 185 characters):
\autocites[§6.1, ¶2]{wald1984}[pp.\,520]{mtw1973}
  [§1, eq.\,4.3]{hawking1974}[§3]{penrose1965}
  [§2]{kerr1963}[ch.\,1]{chandrasekhar1931}

% citemsp equivalent (6 sources, 107 characters):
\citemsp{wald1984/6.1/2, mtw1973/p520,
  hawking1974/1/e4.3, penrose1965/3,
  kerr1963/2, chandrasekhar1931/1}
\end{verbatim}
\end{tcolorbox}

\noindent\textbf{Rendered output:}

\medskip
\noindent
\begin{tabular}{@{}lp{10cm}@{}}
biblatex: & \autocites[§6.1, ¶2]{wald1984}[pp.\,520]{mtw1973}[§1, eq.\,4.3]{hawking1974}[§3]{penrose1965}[§2]{kerr1963}[ch.\,1]{chandrasekhar1931} \\[4pt]
citemsp:  & \citemsp{wald1984/6.1/2, mtw1973/p520, hawking1974/1/e4.3, penrose1965/3, kerr1963/2, chandrasekhar1931/1}
\end{tabular}

\medskip
\noindent At six entries, the \texttt{\textbackslash autocites} source is
78 characters longer (a 42\% reduction) and
produces a bracket sequence that occupies significant horizontal space in the
running text. The \texttt{citemsp} version remains a single bracket with
compact locator annotations.

\subsection*{Summary of advantages}

Three properties distinguish \texttt{\textbackslash citemsp}: (1)~a compact,
uniform syntax that stays readable as entries accumulate; (2)~typed locator
prefixes that visually distinguish sections, equations, pages, theorems, and
other reference types; and (3)~superscript/subscript rendering that keeps
locator information close to each citation number rather than appended as
flat text inside the brackets.


\section{Conventions}

\begin{itemize}
  \item \textbf{Section numbering.} Use the source's native section numbering format (e.g., \texttt{6.1} for ``Section~6.1'').

  \item \textbf{Paragraph counting.} Paragraphs are counted manually from the start of the section. This is inherently edition-dependent and best suited for stable references such as textbooks. Since paragraphs are not explicitly numbered in published sources, these locators may be ambiguous across editions or between readers. When using \P\ locators, consider providing additional context in the surrounding text so the reader can locate the passage independently.

  \item \textbf{Custom prefixes.} Register new locator types with \verb|\citemspprefix{letter}{label}|. The letter must be a single character not already in use. Custom and built-in prefixes can be freely mixed in the same \texttt{\textbackslash citemsp} call.
\end{itemize}


\printbibliography

\end{document}
